\documentclass[a4paper,11pt]{article}
% vitrine_preamble.tex — shared house style for Vitrine compact institutional reports.
% \input this after \documentclass[a4paper,11pt]{article}. Builds with pdflatex/tectonic.
% House palette inherited from report/v5/main.tex (UoSRed = University of Salford red).

\usepackage[utf8]{inputenc}
\usepackage[T1]{fontenc}
\usepackage[british]{babel}
\usepackage{microtype}
\usepackage[a4paper,top=22mm,bottom=22mm,inner=20mm,outer=20mm,headheight=26pt]{geometry}
\usepackage{lmodern}
\usepackage[dvipsnames,svgnames,x11names]{xcolor}

\definecolor{UoSRed}{HTML}{B71234}
\definecolor{DreamLabAccent}{HTML}{6366F1}
\definecolor{DarkBg}{HTML}{0F172A}
\definecolor{MidGrey}{HTML}{64748B}
\definecolor{LightGrey}{HTML}{F1F5F9}
\definecolor{CodeBg}{HTML}{1E293B}
\definecolor{SuccessGreen}{HTML}{15803D}
\definecolor{WarningAmber}{HTML}{B45309}
\definecolor{DangerRed}{HTML}{B91C1C}

\usepackage{graphicx}
\usepackage{booktabs}
\usepackage{tabularx}
\usepackage{array}
\newcolumntype{L}[1]{>{\raggedright\arraybackslash}p{#1}}
\newcolumntype{C}[1]{>{\centering\arraybackslash}p{#1}}
\usepackage{amsmath,amssymb}
\usepackage{csquotes}
\usepackage{enumitem}
\setlist{nosep,leftmargin=1.2em,topsep=2pt}
\usepackage{caption}
\captionsetup{font=small,labelfont={bf,color=UoSRed}}
\usepackage[most]{tcolorbox}
\usepackage{titlesec}
\usepackage{fancyhdr}
\usepackage[colorlinks=true,linkcolor=UoSRed,citecolor=DreamLabAccent,urlcolor=DreamLabAccent,
  breaklinks=true]{hyperref}
\usepackage{enumitem}

% --- Section styling: red rule + tight spacing (compact) ---
\titleformat{\section}{\color{UoSRed}\Large\bfseries}{\thesection}{0.6em}{}[{\color{UoSRed!30}\titlerule[1pt]}]
\titleformat{\subsection}{\color{DarkBg}\large\bfseries}{\thesubsection}{0.5em}{}
\titlespacing*{\section}{0pt}{10pt}{5pt}
\titlespacing*{\subsection}{0pt}{7pt}{3pt}

% --- Header / footer ---
\pagestyle{fancy}\fancyhf{}
\renewcommand{\headrulewidth}{0.6pt}
\renewcommand{\headrule}{\hbox to\headwidth{\color{UoSRed}\leaders\hrule height \headrulewidth\hfill}}
\fancyhead[L]{\small\color{MidGrey}\VitrineDocKicker}
\fancyhead[R]{\small\color{MidGrey}\textbf{\color{UoSRed}Vitrine}}
\fancyfoot[L]{\scriptsize\color{MidGrey}DreamLab AI Consulting Ltd \textperiodcentered{} University of Salford}
\fancyfoot[R]{\scriptsize\color{MidGrey}\thepage}
\newcommand{\VitrineDocKicker}{Vitrine}

% --- Callout boxes ---
\newtcolorbox{keyfinding}[1][]{colback=UoSRed!5,colframe=UoSRed,boxrule=0.8pt,arc=2pt,
  left=6pt,right=6pt,top=4pt,bottom=4pt,fonttitle=\bfseries\color{white},
  coltitle=white,colbacktitle=UoSRed,title={#1}}
\newtcolorbox{technote}[1][]{colback=DreamLabAccent!6,colframe=DreamLabAccent,boxrule=0.7pt,arc=2pt,
  left=6pt,right=6pt,top=4pt,bottom=4pt,fonttitle=\bfseries\color{white},
  coltitle=white,colbacktitle=DreamLabAccent,title={#1}}
\newtcolorbox{riskbox}[1][]{colback=WarningAmber!7,colframe=WarningAmber,boxrule=0.7pt,arc=2pt,
  left=6pt,right=6pt,top=4pt,bottom=4pt,fonttitle=\bfseries\color{white},
  coltitle=white,colbacktitle=WarningAmber,title={#1}}

% --- Title block macro: \vitrinetitle{Title}{Subtitle}{Audience line} ---
\newcommand{\vitrinetitle}[3]{%
  \begin{tcolorbox}[colback=DarkBg,colframe=DarkBg,arc=3pt,left=14pt,right=14pt,top=12pt,bottom=12pt]
    {\color{white}\Huge\bfseries #1}\\[4pt]
    {\color{LightGrey}\large #2}\\[8pt]
    {\color{DreamLabAccent}\small\bfseries #3}
  \end{tcolorbox}\vspace{4pt}%
}

\newcommand{\code}[1]{{\small\ttfamily #1}}
\setlength{\parskip}{4pt}\setlength{\parindent}{0pt}

\renewcommand{\VitrineDocKicker}{Digital IT briefing}

\begin{document}

\vitrinetitle{Vitrine}%
{Capture-adaptive photo/video to structured, textured, game-ready 3D scenes for cultural-heritage digitisation}%
{Briefing for University of Salford Digital IT \textbullet\ Deployment, security posture and platform justification \textbullet\ 2026-07-02}

\begin{keyfinding}[Executive summary]
Vitrine turns ordinary phone photos and video into structured, textured, game-ready 3D
scenes for Unreal Engine 5.8 and Blender, for on-premises cultural-heritage digitisation.
It ships as a \textbf{single hardened Docker image} with \textbf{no SaaS dependency}: the web
UI binds \code{127.0.0.1:7860} and is reached only over an SSH tunnel, so captures, splats
and meshes never leave institution hardware. An internal agent controller runs
\emph{risk-managed} continuous improvement (propose $\to$ adversarially verify $\to$ commit,
with a human owning rebuilds). End-to-end validated 2026-07-02 on 55 hand-held phone frames.
Honest maturity: pre-product, single-site R\&D; capture quality is the dominant quality
bottleneck; per-object segmentation is a known follow-up.
\end{keyfinding}

\section{What Vitrine is and what it produces}

Vitrine is a capture-adaptive pipeline: it diagnoses each input, routes it through
photogrammetry, 3D Gaussian Splatting (3DGS) and mesh-extraction stages as the capture
quality allows, and delivers a scene as \emph{standard game assets} (baked-PBR FBX/GLB, not a
raw splat or point cloud) with explicit \code{v2g:*} provenance metadata. It vendors the
native C++/CUDA 3DGS trainer LichtFeld Studio at a pinned tag (v0.5.3, never forked). The
run below is the validated end-to-end path of 2026-07-02: a brass vessel and ketchup-bottle
still life shot on a phone.

\begin{center}\small
\begin{tabularx}{\textwidth}{@{}L{3.3cm} L{5.0cm} X@{}}
\toprule
\textbf{Stage} & \textbf{Tool / method} & \textbf{Validated output (2026-07-02)} \\
\midrule
Ingest / decode      & DNG/HEIC decode, camera white balance        & 55/55 hand-held phone frames decoded \\
SfM registration     & COLMAP (ALIKED + LightGlue)                   & 100\% frame registration, 10.4k points \\
3DGS training        & LichtFeld Studio v0.5.3 \code{igs+}, 30k iter & GPU-bound $\sim$15 min, 4M-Gaussian room splat \\
Object segmentation  & SAM silhouette crop                           & per-object crop (follow-up area) \\
Object reconstruction& TRELLIS.2 image $\to$ 3D                       & 4096px PBR-textured GLB, 274k verts \\
Scene assembly       & Blender exhibit scene                         & staged scene carrying \code{v2g:*} lineage \\
Web delivery         & \code{<model-viewer>}                          & in-browser preview \\
Engine delivery (opt)& FBX game asset, UE 5.8 (Nanite)               & textured mesh import, no splat runtime \\
\bottomrule
\end{tabularx}
\end{center}

Delivering meshes rather than splats is deliberate: 3DGS in a game engine is a renderer, not
an asset, and carries documented limits (no Lumen GI/reflection interaction, view-dependent
appearance baked in, $\sim$180\,MB per 1M splats vs 5--20\,MB for a textured mesh, no standard
PBR or collision) [13]. Vitrine treats the splat as an intermediate and ships a mesh.

\section{The single hardened mono-image and its attack surface}

For IT signoff the relevant fact is the small, auditable surface. Everything runs inside one
image; there are no ambient external calls, no per-asset cloud storage, and nothing published
to the LAN by default. The design is recorded in ADR-022 (secure single-image) and ADR-023
(web consolidation).

\begin{center}\small
\begin{tabularx}{\textwidth}{@{}L{4.1cm} X@{}}
\toprule
\textbf{Concern} & \textbf{Vitrine posture} \\
\midrule
Deployment unit      & One auditable Docker mono-image; no service sprawl to inventory or patch piecemeal. \\
Network exposure     & Web UI binds \code{127.0.0.1:7860}; host publish pinned to \code{127.0.0.1}. The LAN never sees the service. \\
Remote access        & Operator tunnels in: \code{ssh -N -L 7860:localhost:7860}. No listener on a routable interface. \\
Third-party tooling  & Least-privilege internal venv + OS-user isolation; bundled model/tool code cannot read host or pipeline secrets. \\
Data egress          & None by design. Captures, splats and meshes stay on institution hardware; air-gappable. \\
Provenance / audit   & Every artifact carries \code{v2g:*} lineage; runs are reproducible from pinned SOTA versions. \\
Design records       & ADR-022 (secure single-image), ADR-023 (web consolidation). \\
\bottomrule
\end{tabularx}
\end{center}

\begin{technote}[Why this is IT-signable and air-gappable]
The default posture is \textbf{deny-by-network}: the only way to reach the UI is an
authenticated SSH session that the institution already governs, then a local-forward to
loopback. No inbound port is open to the LAN or internet, so the pipeline can run fully
air-gapped once model weights are staged on disk. Because it is one image with an internal
least-privilege venv and OS-user split, the audit question is bounded: one image to scan, one
loopback binding to confirm, one provenance scheme to inspect. This directly answers the
data-sovereignty and CARE/FAIR requirements that heritage collections impose (\S5).
\end{technote}

\section{Internal agentic resilience: risk-managed continuous improvement}

Vitrine embeds an agent controller that diagnoses each capture, selects the pipeline, drives
the tools, runs in-flight evaluations, and recovers from failure (restart-resilient model
lifecycles, phased VRAM allocation). Improvement is not ad-hoc: the controller fans out
sub-agent ``meshes'' that \textbf{propose} changes and \textbf{adversarially verify} them
before anything is committed --- \textbf{propose $\to$ verify $\to$ commit}, with a human owning
the rebuild. Two concrete cases:

\begin{itemize}
\item \textbf{QE fleet caught a shipping bug pre-commit.} A quality-engineering sub-agent
  fleet found a dead API-blueprint wiring bug during adversarial verification and blocked it
  before it reached a build.
\item \textbf{Self-healing container entrypoint.} When a pinned wheel's ABI drifted under a
  Torch upgrade, the entrypoint rebuilt the affected CUDA extension from source rather than
  failing the run, keeping the image reproducible.
\end{itemize}

This is automated but bounded: agents propose and verify; a human authorises the rebuild that
ships. It is a maintenance and reliability property, not autonomy over production.

\begin{riskbox}[Residual risks and mitigations]
\textbf{Risk:} agentic self-modification could regress the pipeline. \textbf{Mitigation:} no
change ships without adversarial verification and a human-owned rebuild; nothing auto-deploys.
\quad\textbf{Risk:} dependency/ABI drift breaks a run. \textbf{Mitigation:} exact-pinned SOTA
versions plus self-healing rebuild-from-source in the entrypoint. \quad\textbf{Risk:} output
quality. \textbf{Mitigation (and honest limit):} quality is capture-limited, not
pipeline-limited --- the dominant lever is capture, and per-object silhouette segmentation is a
known, tracked follow-up. Vitrine is single-site R\&D, not a shipping product.
\end{riskbox}

\section{Position versus competitors}

No mainstream tool ships the whole capture-to-game-asset path as one sovereign, offline unit.
Cloud tools are structurally disqualified for sovereign heritage data; strong desktop tools
are single apps, not an integrated pipeline; the closest academic peer stops at splats.

\begin{center}\small
\begin{tabularx}{\textwidth}{@{}L{3.0cm} C{1.9cm} C{1.9cm} C{1.9cm} C{2.1cm} C{1.6cm}@{}}
\toprule
\textbf{Tool / service} & \textbf{On-prem?} & \textbf{Trains 3DGS?} & \textbf{Per-object?} & \textbf{UE game asset?} & \textbf{Sovereign?} \\
\midrule
\textbf{Vitrine}          & Yes         & Yes         & Yes (follow-up) & Yes (FBX/Nanite) & Yes \\
Luma AI                   & No (cloud)  & Yes         & No          & Splat only     & No \\
Polycam                   & No (cloud)  & Yes         & No          & Mesh, no scene & No \\
KIRI Engine               & No (cloud)  & Yes         & Object mesh & No             & No \\
Scaniverse (Niantic)      & On-device   & Yes         & No          & No             & Partial \\
RealityScan (Epic)        & Yes         & No (SfM)    & No          & Mesh, manual   & Yes \\
Agisoft Metashape         & Yes         & No (export) & No          & Mesh, manual   & Yes \\
Postshot (Jawset)         & Yes         & Yes         & No          & Splat only     & Yes \\
Gaussian Heritage [11]    & Yes (Docker)& Yes         & Yes         & Splat only     & Yes \\
\bottomrule
\end{tabularx}
\end{center}

The moat is the \emph{combination}, not any single axis: sovereign on-prem delivery, 3DGS
training, per-object scene structure, and UE 5.8 game-asset output in one deployable image.
The nearest neighbours are RealityScan (on-prem + UE, but no 3DGS training, no automated
object scene) [12] and Gaussian Heritage (structuring + Docker, but research-grade splats, no
game-asset delivery) [11]; no competitor spans both. Where rivals are genuinely ahead ---
stated plainly --- is capture UX and ergonomics, out-of-box per-object mesh quality, 3DGS
trainer polish, and scale/track record. Vitrine's object quality is still capture-limited and
iterating.

\section{Business case: why on-prem sovereign digitisation}

Two decay vectors make structured, open, re-renderable 3D surrogates a preservation asset
rather than a novelty.

\textbf{Physical decay.} UNESCO/WRI (July 2025) find \textbf{73\% of 1,172 non-marine World
Heritage sites are highly exposed to water-related hazards, and 21\% face multiple overlapping
risks} [1]. Catastrophic single-event losses recur: the National Museum of Brazil fire (2018)
destroyed most of a $\sim$20-million-item collection, recoverable only through prior imagery
[3]; after Notre-Dame (2019), Andrew Tallon's 2010 laser scan ($>$1 billion points,
$\sim$5\,mm) became the restoration blueprint [2]. The lesson is consistent: capture beats
reconstruction, and the window in which an object still exists and has not yet been captured
is closing.

\textbf{Digital bitrot.} Digitising is necessary but not sufficient. Pew (2024) measured that
\textbf{38\% of web pages that existed in 2013 were gone by 2023} [4]; the Digital Preservation
Coalition's 2023 ``Bit List'' reached \textbf{87 endangered digital categories} and reclassified
unpublished research data as \emph{Critically Endangered} [5]; storage media are not archival
(hard drives 3--5 years, LTO tape 10--20 in practice) so only redundancy plus active migration
survives [6]. Open, standardised formats age far better --- glTF is an ISO/IEC 12113:2022
standard and OpenUSD is vendor-neutral [10]. Vitrine emits a structured, format-migratable
asset stack (mesh + baked textures + \code{v2g:*} lineage) that a future institution can still
open; a defunct SaaS account is not.

\textbf{Why sovereignty, not cloud.} Roughly \textbf{5--10\% of collections are on display};
the rest sit in storage (90--99\% by museum type) [7], so any serious programme is
high-volume, and recurring per-asset cloud fees scale badly against it. Sacred,
rights-encumbered or community-owned material cannot be governed under a third party's terms:
CARE principles and Indigenous-data-sovereignty practice require the originating community to
control the infrastructure and access, not merely receive a file on someone else's server [9].
FAIR/paradata scholarship treats the digitisation \emph{process} as a research object needing
documented provenance --- which Vitrine's pinned, containerised, \code{v2g:*}-tagged runs
provide. Demand is documented: the University of Glasgow's \pounds5.6m ``Museums in the
Metaverse'' surveyed 2,000+ people and found \textbf{79\% interested in digital tools to
explore collections not on display} [8].

\section{Funding horizon}

Vitrine's four-way hook --- heritage digitisation, immersive/games-engine delivery (UE 5.8,
Nanite), sovereign/on-prem AI, and real games-asset workflow --- maps onto a realistic ladder
of live opportunities. UK entities are associated to Horizon Europe for 2024--2027 and can bid
Cluster 2 directly.

\begin{center}\small
\begin{tabularx}{\textwidth}{@{}L{5.4cm} L{3.0cm} X@{}}
\toprule
\textbf{Programme / call} & \textbf{Funder} & \textbf{Fit / positioning} \\
\midrule
HORIZON-CL2-2027-01-HERITAGE-08 (intangible CH; $\sim$Sep 2027) & Horizon Europe CL2 [14] & Primary anchor: Vitrine is the technical pipeline WP inside a culture-led consortium. \\
HORIZON-CL2-2026-01-HERITAGE-03 (AI in CCSI practice; $\sim$Sep 2026) & Horizon Europe CL2 [14] & Cleaner ``efficient AI workflow'' fit, one year sooner; data-sovereign case study. \\
N-RICH Prototype / Towards a National Collection (2025--27) & UKRI / AHRC [15] & Sovereign, cyber-hardened on-prem digitisation capability answers N-RICH frameworks. \\
CoSTAR Network (Live/National/Realtime Labs) & UKRI / AHRC [16] & Supplies game-ready structured assets for real-time / virtual-production pipelines. \\
Innovate UK Creative Catalyst (\pounds150--200k) & Innovate UK & DreamLab AI leads, UoS collaborates; productise the capture-to-UE tool. \\
MITIH (MediaCity Immersive Tech Hub) & Innovate UK / UoS & Immediate, local, low-friction proof point on a real GM collection. \\
National Lottery Heritage Grants (rolling, $\le$\pounds250k) & NLHF & Fund a cultural partner to run a real digitisation case study. \\
\bottomrule
\end{tabularx}
\end{center}

\textbf{Positioning.} Sequence the domestic funders as rungs that de-risk the 2027 EU anchor:
MITIH and Innovate UK for near-term productisation and a local proof point; CoSTAR and AHRC/N-RICH
to embed Vitrine as shared real-time capability and a national sovereign-digitisation asset;
National Lottery Heritage Grants to generate the real-world case histories the Community-of-Practice
pilot bid promised. Each rung produces the partnerships and evidence the HERITAGE-08 consortium bid
needs, and every rung leads with the same sovereign, own-your-data framing that the hardened
mono-image makes credible.

\section*{References}
\small
\begin{enumerate}[leftmargin=2.2em,label={[\arabic*]}]
\item UNESCO/WRI --- \emph{Nearly three-quarters of World Heritage sites at high risk from water-related hazards} (2025). \url{https://whc.unesco.org/en/news/2788}
\item Tallon Notre-Dame laser scan (2010; $>$1B points, $\sim$5\,mm; restoration blueprint) --- GIM International. \url{https://www.gim-international.com/content/news/the-indispensable-role-of-point-cloud-data-in-rebuilding-notre-dame}
\item National Museum of Brazil fire (collection scale; 3D reconstruction from prior imagery). \url{https://en.wikipedia.org/wiki/National_Museum_of_Brazil_fire}
\item Pew Research --- \emph{When Online Content Disappears} (link rot, 2024). \url{https://www.pewresearch.org/data-labs/2024/05/17/when-online-content-disappears/}
\item Digital Preservation Coalition --- \emph{Global Bit List 2023} (87 endangered categories). \url{https://www.dpconline.org/our-work/digitally-endangered-species}
\item CLIR Pub 54 --- \emph{Life Expectancy: How Long Will Magnetic Media Last?} \url{https://www.clir.org/pubs/reports/pub54/4life_expectancy/}
\item ICOM / ICCROM --- \emph{Collections in Storage} ($\sim$5--10\% displayed). \url{https://icom.museum/wp-content/uploads/2024/05/Report_ICOM-STORAGE_EN_Final.pdf}
\item University of Glasgow --- \emph{Museums in the Metaverse} (\pounds5.6m; 79\% interested), via Blooloop. \url{https://blooloop.com/museum/in-depth/emerging-museum-technologies/}
\item Indigenous data sovereignty / CARE and digital repatriation. \url{https://terentia.io/blog/indigenous-data-sovereignty-cultural-heritage} ; \url{https://en.wikipedia.org/wiki/Digital_repatriation}
\item Khronos --- OpenUSD/glTF interoperability (glTF = ISO/IEC 12113:2022). \url{https://www.khronos.org/blog/building-bridges-in-3d-aousd-and-khronos-collaborate-on-openusd-and-gltf-interoperability}
\item Gaussian Heritage (3DGS + integrated per-object segmentation; Docker), arXiv 2409.19039. \url{https://arxiv.org/abs/2409.19039}
\item RealityScan 2.1 (on-prem photogrammetry; COLMAP/RF export; does not train 3DGS) --- CG Channel. \url{https://www.cgchannel.com/2025/11/epic-games-releases-realityscan-2-1/}
\item 3DGS-in-engine limitations and the mesh path. \url{https://www.polyvia3d.com/guides/gaussian-splatting-unity-unreal} ; \url{https://www.strayspark.studio/blog/gaussian-splatting-unreal-engine-5-capture-to-game-pipeline}
\item Horizon Europe Work Programme 2026--2027, Cluster 2 (HERITAGE topics). \url{https://ec.europa.eu/info/funding-tenders/opportunities/docs/2021-2027/horizon/wp-call/2026-2027/wp-5-culture-creativity-and-inclusive-society_horizon-2026-2027_en.pdf}
\item N-RICH Prototype / Towards a National Collection (UKRI/AHRC). \url{https://www.nationalcollection.org.uk/n-rich-prototype}
\item CoSTAR --- Convergent Screen Technologies and Performance in Realtime (UKRI/AHRC). \url{https://www.ukri.org/councils/ahrc/remit-programmes-and-priorities/convergent-screen-technologies-and-performance-in-realtime-costar/}
\end{enumerate}

\end{document}
